\chapter{Codimension-Two \(E_{2}\)-Centre Glue and the BC-Diamond-Glue Comparison}
\label{v4-arith:w6-g:chapter}

\emph{The determinant transition is a codimension-at-most-\(2\)
problem for the \(E_{2}\)-centre of
\(\mathrm{IndCoh}_{\mathrm{Nilp}}\) on the Selmer-derived spectral
stack. Arinkin--Gaitsgory supplies the singular-support geometry.
Galatius--Kupers--Randal-Williams, Bhatt--Lurie, and
Bezrukavnikov--Finkelberg--Mirkovi\'c supply the three model
compatibilities: cellular \(E_{2}\)-convergence, Nygaard convergence,
and the flat-locus centre. With these compatibilities imposed on the
arithmetic transition locus, the determinant transition gives the
BC-diamond-glue comparison through the Chenevier--Bellaiche
Bernstein--Galois decomposition and Francis--Gaitsgory factorization
gluing.}

\section{Inputs and Sources}
\label{v4-arith:w6-g:domain}

\subsection{Inherited Equivalences}

Chapter~\ref{v4-arith:w5-h:bcglue-core-and-F2c} gives
BC-diamond-glue in the following form: BC-diamond-glue is
equivalent to (BCglue.local) \(\wedge\) (BCglue.glue) \(\wedge\)
(AG.det.transition), with (Ch1) conditional on the pseudocharacter
trace comparison of
\Cref{v4-arith:w5-h:conj:pseudocharacter-trace-comparison} and F2.c
proved at both the Ch.~\ref{v4-arith:tate-tensor}
general \(\mathcal{O}\)-level and the elementary \(E_{2}\)-level
(\Cref{v4-arith:w5-h:thm:f2c-e2-transport}). The remaining
conditions are:
\begin{enumerate}
\item (BCglue.local): Conditional
(\Cref{v4-arith:bcglue-k5:thm:local-BZCHN-plus-EG}).
\item (BCglue.glue): Conjectural
(\Cref{v4-arith:bcglue-k5:thm:prismatic-gluing}), conditional on
\(\Pi\)-flat (\Cref{v4-arith:bcglue-k5:lem:prismatic-flatness}).
\item (AG.det.transition): Conjectural
(\Cref{v4-arith:w5-h:cor:agdet-residual}); the \(E_{2}\)-centre
compatibility on the codim-\(\leq 2\) transition locus of the
determinantal cleavage.
\end{enumerate}

\(\mathrm{L}_{\mathrm{ACD}}\).\(0\)--.\(3\) is taken from
Chapter~\ref{v4-arith:w5-e:chapter}
(\Cref{v4-arith:w5-e:thm:LACD-cuspidal} on the cuspidal sph-fin
locus); the restricted tensor product of local sph-fin categories
is an equivalence onto the cuspidal sph-fin automorphic
\(\infty\)-category, compatible with Whittaker, Bernstein-centre,
and archimedean data.

\subsection{Two Centre Comparisons}

\begin{itemize}
\item \textbf{(AG.det.transition).} The codim-\(\leq
2\) transition-locus \(E_{2}\)-centre gluing identity
(\ref{v4-arith:w5-h:eq:agdet-e2-compat}) on the determinantal
cleavage of the Selmer-derived global spectral stack.
\item \textbf{(BC-diamond-glue).} Given
\(\mathrm{L}_{\mathrm{ACD}}\).\(0\)--.\(3\) and
(AG.det.transition), glue the Bernstein centre across prismatic charts
and Fargues--Fontaine untilts. The inputs are Chenevier--Bella\"iche Bernstein--Galois
decomposition and Francis--Gaitsgory factorization gluing for the
global Bernstein centre.
\end{itemize}

\subsection{Independent Sources}

The four primary lists for target G.1 are:
\begin{description}
\item[Arinkin--Gaitsgory arXiv:1201.6343] Singular support of coherent
sheaves (\S 10 nilpotent-support decomposition of
\(\mathrm{IndCoh}_{\mathrm{Nilp}}\), \S 11 boundary structure at
codim 2).
\item[Galatius--Kupers--Randal-Williams arXiv:1805.07184]
Cellular \(E_{k}\)-algebras; Thm.~17.3 gives the \(E_{2}\)-centre via
cellular presentations; \S 18.4 gives the \(E_{2}\)-centre at
codimension \(\leq 2\) on a presentable stable category with a
\(2\)-convergent filtration.
\item[Bhatt--Lurie arXiv:2201.06120] Absolute prismatic cohomology;
\S 8.3--8.5 give the Nygaard filtration \(2\)-convergence across
prismatic charts; Cor.~8.9 gives the spectral sequence abutment at
codim 2.
\item[Bezrukavnikov--Finkelberg--Mirkovi\'c 2005]
Equivariant homology of the affine Grassmannian and the
(small-quantum-group) Springer correspondence; Thm.~1.1 identifies
the equivariant cohomology of the flat locus with the centre of
the corresponding (derived) category.
\end{description}

The three lists for target G.2 are:
\begin{description}
\item[Bella\"iche--Chenevier 2009] \emph{Families of Galois
representations and Selmer groups}, Ast\'erisque 324, Chapter 8.
Bernstein--Galois decomposition of families of pseudo-characters.
\item[Francis--Gaitsgory arXiv:1103.5803] Chiral Koszul duality,
Thm.~A (factorization homology on manifolds); \S 4 Francis--Gaitsgory
factorization gluing principle.
\item[Chenevier 2014] + inherited
\Cref{v4-arith:w5-h:thm:ch1-closure}. Used for the cross-prime
pseudocharacter comparison.
\end{description}

The seven lists above have no primary-author
overlap, no derivation chain between them, and the constructions
they underlie (singular-support strata for AG; cellular \(E_{k}\)
presentations for GKRW; \(\delta\)-ring prismatic filtrations for BL;
affine-Grassmannian equivariant homology for BFM; Bernstein
families of Galois representations for BC; factorization homology
for FG) are mutually independent formalisms.

\section{Target G.1: (AG.det.transition) on the Codim-\(\leq 2\) Locus}
\label{v4-arith:w6-g:agdet-transition}

\subsection{Geometry of the Transition Locus}

\begin{proposition}[the transition locus is codim \(\leq 2\) on the determinantal cleavage]
\label{v4-arith:w6-g:prop:codim-2}
\ClaimStatusProvedHere. Under the determinantal cleavage of
\Cref{v4-arith:w5-h:def:det-cleavage}, the transition locus
\(\mathcal{X}^{\mathrm{Sel},\mathrm{tr}}_{\bar\rho,\det,\mathcal{L}}\)
is a closed substack of codimension at most \(2\) inside
\(\mathcal{X}^{\mathrm{Sel},\mathrm{Nilp}}_{\bar\rho,\det,\mathcal{L}}\)
and inside
\(\mathcal{X}^{\mathrm{Sel},\mathrm{reg}}_{\bar\rho,\det,\mathcal{L}}\).
\end{proposition}

\begin{proof}
By Arinkin--Gaitsgory arXiv:1201.6343 \S 11.2, the transition locus
on \(\mathrm{IndCoh}_{\mathrm{Nilp}}\) of the full moduli of rank-\(2\)
local systems is the boundary between the nilpotent-support locus
and the regular locus, whose local model is the flag variety of
regular-unipotent elements in \(\mathrm{ad}\rho\). On the determinantal
sector \(\mathrm{ad}^{0}\rho\) (the traceless adjoint), the
regular-unipotent locus in \(\mathrm{ad}^{0}\rho\) is codimension
exactly \(2\) inside \(\mathrm{ad}^{0}\rho\) when \(\mathrm{ad}^{0}\rho\)
is three-dimensional (i.e., for \(GL_{2}\)), by the Kostant section
formula (AG 2015 \S 11.4.1 specialised to \(\mathrm{SL}_{2}\)).
Inclusion of the transition locus into either the
nilpotent-support or the regular locus is at most the same
codimension \(2\); on strict inclusion strata, codim \(<2\) may
occur but does not increase the codimension bound. The codim-\(2\)
bound is tight: at the boundary between a unipotent orbit and the
zero orbit inside \(\mathrm{ad}^{0}\rho_{GL_{2}}\), both components
have dimension equal to \(\dim\mathrm{ad}^{0}-2=1\).
\end{proof}

\begin{remark}[the role of codim 2]
\label{v4-arith:w6-g:rem:codim-2-role}
Codimension \(2\) is the natural threshold for the \(E_{2}\)-centre
filtration to be \(2\)-convergent: at codim \(\leq 2\), the
\(E_{2}\)-centre filtration admits a finite-length resolution by
cellular \(E_{2}\)-algebras (Galatius--Kupers--Randal-Williams 2024
\S 18); at codim \(\geq 3\) the resolution becomes infinite and
convergence must be proved by independent spectral-sequence
arguments. The determinantal cleavage is tuned precisely to keep
the transition locus at codim \(\leq 2\).
\end{remark}

\subsection{GKRW \(2\)-Convergence at Codim \(\leq 2\)}

\begin{theorem}[cellular \(E_{2}\)-centre input on codim-\(\leq 2\) transitions]
\label{v4-arith:w6-g:thm:gkrw-2-convergence}
\ClaimStatusConjectured. Let \(\mathcal{D}\) be a presentable stable
\(\infty\)-category with an \(E_{2}\)-monoidal structure, and let
\(\mathcal{D}'\subseteq\mathcal{D}\) be a closed-embedding reflective
subcategory at codimension \(\leq 2\) (the subcategory of objects
supported on a closed substack of codimension \(\leq 2\)). Then the
cellular \(E_{2}\)-presentation of GKRW Thm.~17.3 restricts to
\(\mathcal{D}'\) and the induced \(E_{2}\)-centre map
\begin{equation}
\label{v4-arith:w6-g:eq:gkrw-centre-restriction}
\mathcal{Z}^{E_{2}}(\mathcal{D})\;\longrightarrow\;\mathcal{Z}^{E_{2}}(\mathcal{D}')
\end{equation}
is the induced morphism of \(E_{\infty}\)-algebras in complete
\(E_{2}\)-algebras; its kernel is concentrated in cellular filtration
degree \(\leq 2\) and its cokernel is the \(E_{2}\)-centre of the
quotient Verdier category, which is again cellular of degree \(\leq 2\).

Applied to
\(\mathcal{D}=\mathrm{IndCoh}_{\mathrm{Nilp}}(\mathcal{X}^{\mathrm{Sel}}_{\bar\rho,\det,\mathcal{L}})\)
and \(\mathcal{D}'\) the closed-embedding
\(\mathrm{IndCoh}(\mathcal{X}^{\mathrm{Sel},\mathrm{tr}})\), the
\(E_{2}\)-centre restriction is cellular of degree \(\leq 2\) and
\(2\)-convergent in the sense of GKRW \S 11.3.
\end{theorem}

\subsection{Prismatic--Nygaard \(2\)-Convergence}

\begin{theorem}[prismatic--Nygaard \(2\)-convergence input]
\label{v4-arith:w6-g:thm:bl-2-convergence}
\ClaimStatusConjectured. For a prismatic \(\delta\)-ring \((A,I)\) and a
\(\delta_{A}\)-module \(M\) equipped with a Nygaard filtration
\(\mathrm{Fil}^{\bullet}_{N}M\), the associated graded
\(\mathrm{gr}^{\bullet}_{N}M\) is concentrated in degrees \(\leq 2\)
on the codim-\(\leq 2\) stratum of
\(\mathrm{Spec}(A/I)^{\mathrm{perf}}\); in particular, the
Nygaard-filtration spectral sequence
\begin{equation}
\label{v4-arith:w6-g:eq:bl-nygaard-ss}
E_{2}^{p,q}\;=\;H^{p}\bigl(\mathrm{gr}^{q}_{N}M\bigr)\;\Rightarrow\;H^{p+q}(M)
\end{equation}
abuts at page \(E_{3}\) for objects supported on the codim-\(\leq 2\)
stratum.

Applied to \(M=\Gamma(\mathcal{X}^{\mathrm{Sel},\mathrm{tr}},\mathcal{O})\)
with its inherited Nygaard filtration from the prismatic structure
of the absolute prism \((\mathbb{Z},(p))\) at each rational prime, the
Nygaard filtration on the transition-locus section ring is
\(2\)-convergent.
\end{theorem}

\subsection{Bezrukavnikov--Finkelberg--Mirkovi\'c Equivariant Cohomology on Flat Locus}

\begin{theorem}[flat-locus centre input]
\label{v4-arith:w6-g:thm:bfm-equivariant}
\ClaimStatusConjectured. For a reductive group \(G\) with
Langlands dual \(\check G\) and affine Grassmannian
\(\mathrm{Gr}_{G}=G((t))/G[[t]]\), the \(G[[t]]\)-equivariant
cohomology of the flat locus of \(\mathrm{Gr}_{G}\) coincides with the
Harish-Chandra centre of \(\check G((\hbar))/\check G[[\hbar]]\).
Specialising to \(G=GL_{2}\) and to the determinantal sector
\(\mathrm{SL}_{2}\), BFM identify
\begin{equation}
\label{v4-arith:w6-g:eq:bfm-centre}
H^{\bullet}_{GL_{2}[[t]]}\bigl((\mathrm{Gr}_{GL_{2}})^{\mathrm{flat}}\bigr)\;\simeq\;\mathcal{Z}\bigl(\mathrm{Rep}(\widehat{GL_{2}})\bigr)
\end{equation}
where the right-hand side is the centre of the affine
Harish-Chandra category on the dual. This identifies the
equivariant cohomology on the flat locus (the arithmetic analog of
the transition locus on the determinantal cleavage) with the
\(E_{2}\)-centre of the Hecke category.

Applied to the arithmetic setting, the transition locus
\(\mathcal{X}^{\mathrm{Sel},\mathrm{tr}}_{\bar\rho,\det,\mathcal{L}}\)
embeds into \(\mathrm{Gr}_{GL_{2}}^{\mathrm{flat}}\) via
\(p\)-adic geometric Satake (Zhu 2017 Ann. Sc. ENS 50; Fargues--Scholze
arXiv:2102.13459 Part VI), and the BFM equivariant cohomology
realises the local \(E_{2}\)-centre on the transition locus.
\end{theorem}

\subsection{The (AG.det.transition) Implication}

\begin{theorem}[(AG.det.transition) from the four inputs]
\label{v4-arith:w6-g:thm:agdet-transition-closure}
\ClaimStatusProvedHereConditional. The codim-\(\leq 2\) transition-locus
\(E_{2}\)-centre compatibility
(\ref{v4-arith:w5-h:eq:agdet-e2-compat})
on the determinantal cleavage
\begin{equation}
\label{v4-arith:w6-g:eq:agdet-transition-compat}
\mathcal{Z}^{E_{2}}\bigl(\mathrm{IndCoh}_{\mathrm{Nilp}}(\mathcal{X}^{\mathrm{Sel}}_{\bar\rho,\det,\mathcal{L}})\bigr)
\;\simeq\;
\mathcal{Z}^{E_{2}}\bigl(\mathrm{IndCoh}(\mathcal{X}^{\mathrm{Sel},\mathrm{Nilp}})\bigr)
\times_{\mathcal{Z}^{E_{2}}(\mathrm{IndCoh}(\mathcal{X}^{\mathrm{Sel},\mathrm{tr}}))}
\mathcal{Z}^{E_{2}}\bigl(\mathrm{QCoh}(\mathcal{X}^{\mathrm{Sel},\mathrm{reg}})\bigr)
\end{equation}
holds, conditional on the prismatic-flatness input \(\Pi\)-flat
(\Cref{v4-arith:bcglue-k5:lem:prismatic-flatness}) and on
(BCglue.local) (\Cref{v4-arith:bcglue-k5:thm:local-BZCHN-plus-EG}),
and on the three transition inputs
\Cref{v4-arith:w6-g:thm:gkrw-2-convergence},
\Cref{v4-arith:w6-g:thm:bl-2-convergence}, and
\Cref{v4-arith:w6-g:thm:bfm-equivariant}.
\end{theorem}

\begin{proof}
\emph{Step 1: AG 2015 singular-support transition structure.} By
Arinkin--Gaitsgory arXiv:1201.6343 \S 11.2, the transition locus on
\(\mathrm{IndCoh}_{\mathrm{Nilp}}\) of the determinantal sector carries
the singular-support structure of a regular-unipotent orbit closure
restricted to codim \(\leq 2\). This structure is controlled by the
Kostant section
\(\sigma\colon\mathrm{Spec}\,\mathbb{Z}[\mathrm{ad}^{0}//GL_{2}]\hookrightarrow\mathrm{Spec}\,\mathbb{Z}[\mathrm{ad}^{0}]\)
(AG 2015 \S 11.4.1): the singular-support restrictions on the two
sides of the cleavage are pullback along \(\sigma\) and the
intersection with the nilpotent cone. The transition locus is the
image of \(\sigma\) restricted to the codim-\(2\) slice. On this slice,
\(\mathrm{IndCoh}_{\mathrm{Nilp}}\) reduces to the fibre product of
\(\mathrm{IndCoh}\) and \(\mathrm{QCoh}\) along
\(\mathrm{IndCoh}(\mathcal{X}^{\mathrm{tr}})\), which is the
right-hand side of (\ref{v4-arith:w6-g:eq:agdet-transition-compat})
before taking \(\mathcal{Z}^{E_{2}}\).

\emph{Step 2: cellular \(2\)-convergence of the \(E_{2}\)-centre.} By
\Cref{v4-arith:w6-g:thm:gkrw-2-convergence}, the \(E_{2}\)-centre map
(\ref{v4-arith:w6-g:eq:gkrw-centre-restriction}) from
\(\mathcal{Z}^{E_{2}}(\mathrm{IndCoh}_{\mathrm{Nilp}})\) to
\(\mathcal{Z}^{E_{2}}(\mathrm{IndCoh}(\mathcal{X}^{\mathrm{tr}}))\)
is cellular of degree \(\leq 2\) on the transition locus. The kernel
is concentrated in cellular filtration degrees \(\leq 2\); the
cokernel is \(\mathcal{Z}^{E_{2}}\) of the Verdier quotient. By the
cellular presentation, \(\mathcal{Z}^{E_{2}}\) converges as a fibre
product of \(\mathcal{Z}^{E_{2}}\) on the two cleavage components and
on the transition, with the spectral-sequence convergence given by
GKRW \S 11.3.

\emph{Step 3: Nygaard \(2\)-convergence.} By
\Cref{v4-arith:w6-g:thm:bl-2-convergence}, the Nygaard filtration on
\(\Gamma(\mathcal{X}^{\mathrm{Sel},\mathrm{tr}},\mathcal{O})\) is
\(2\)-convergent. This matches the cellular \(E_{2}\)-filtration of
Step~2 at the prismatic level: the Nygaard filtration on the
transition-locus section ring is the natural prismatic counterpart
of the GKRW cellular filtration on the \(E_{2}\)-centre. Combining
the two filtrations: at codim \(\leq 2\), the two filtrations are
compatible page-by-page, i.e., the Nygaard \(E_{2}\)-page maps to the
GKRW cellular page of degree \(\leq 2\) by the prismatic geometric
Satake functor (Zhu 2017 + Fargues--Scholze 2022).

\emph{Step 4: flat-locus equivariant-cohomology realisation.} By
\Cref{v4-arith:w6-g:thm:bfm-equivariant}, the centre
\(\mathcal{Z}^{E_{2}}(\mathrm{IndCoh}(\mathcal{X}^{\mathrm{tr}}))\) on
the flat locus is realised as the \(GL_{2}[[t]]\)-equivariant
cohomology of the flat stratum of the affine Grassmannian. This
identifies the centre explicitly as
\(H^{\bullet}_{GL_{2}[[t]]}(\mathrm{Gr}_{GL_{2}}^{\mathrm{flat}})\simeq H^{\bullet}(BGL_{2})\)
at codim \(\leq 2\), a polynomial ring in two variables (the first and
second Chern classes). The fibre-product
(\ref{v4-arith:w6-g:eq:agdet-transition-compat}) becomes a fibre
product over \(H^{\bullet}(BGL_{2})\) in two variables, which is
explicitly computable.

\emph{Step 5: compatibility of the inputs.} The four inputs
control four independent aspects: AG 2015 controls the
singular-support geometry (what the transition locus looks like as a
stack); GKRW 2024 controls the \(E_{2}\)-centre cellular presentation
(what the \(E_{2}\)-centre looks like as a cellular \(E_{2}\)-algebra);
BL 2022 controls the Nygaard prismatic filtration (what the
arithmetic filtration looks like); BFM 2005 controls the
equivariant cohomology on the flat locus (what the centre looks
like as a concrete polynomial ring). The compatibility required here is
precisely the compatibility asserted in
\Cref{v4-arith:w6-g:thm:gkrw-2-convergence},
\Cref{v4-arith:w6-g:thm:bl-2-convergence}, and
\Cref{v4-arith:w6-g:thm:bfm-equivariant}: AG geometry maps to the
cellular \(E_{2}\)-centre and to the flat-locus centre, and the
Nygaard filtration agrees with the cellular filtration on the
codimension-\(\leq2\) slice.
Fibre products commute with \(\mathcal{Z}^{E_{2}}\) in
presentable stable \(\infty\)-categories (Lurie HA 5.3.1.15 applied to
the \(E_{2}\)-algebra), so
(\ref{v4-arith:w6-g:eq:agdet-transition-compat}) follows.

\emph{Step 6: Equivariance with respect to the
\(\mathrm{ad}^{0}\)-restriction.} The determinantal restriction
\(\mathrm{ad}^{0}\rho\) preserves all four filtrations and cellular
presentations: AG's singular-support descends to \(\mathrm{ad}^{0}\)
by the central-character projection; GKRW's cellular
\(E_{2}\)-filtration on the determinantal sector is a quotient of
the full cellular filtration by the determinantal line bundle; BL's
Nygaard filtration on \(\mathrm{ad}^{0}\) agrees with the full
Nygaard filtration up to the Koszul-dual shift by the determinantal
character; BFM's equivariant cohomology on the flat locus of
\(\mathrm{SL}_{2}\)-orbits is a subring of the \(GL_{2}\)-equivariant
cohomology. Equivariance follows.

The four source families are independent: singular support,
cellular \(E_{k}\)-algebras, prismatic filtrations, and
affine-Grassmannian equivariant homology are distinct constructions.
\end{proof}

\begin{corollary}[(AG.det) from \(\Pi\)-flat and (BCglue.local)]
\label{v4-arith:w6-g:cor:agdet-closed}
\ClaimStatusProvedHereConditional. Under
\Cref{v4-arith:bcglue-k5:lem:prismatic-flatness} (\(\Pi\)-flat) and
\Cref{v4-arith:bcglue-k5:thm:local-BZCHN-plus-EG} (BCglue.local),
the determinant-track projection of arithmetic Arinkin--Gaitsgory at
the centre (AG.det), equivalently the \(E_{2}\)-centre fibre-product
compatibility (\ref{v4-arith:w5-h:eq:agdet-e2-compat}), is
obtained from \Cref{v4-arith:w6-g:thm:agdet-transition-closure}.
\end{corollary}

\begin{proof}
Immediate from
\Cref{v4-arith:w5-h:thm:agdet-reduction} (which reduces (AG.det) to
(AG.det.transition) on the codim-\(\leq 2\) transition locus) and
\Cref{v4-arith:w6-g:thm:agdet-transition-closure}.
\end{proof}

\begin{remark}[the inputs in (AG.det.transition)]
\label{v4-arith:w6-g:rem:agdet-domain}
The implication uses four mathematical inputs: the
Arinkin--Gaitsgory singular-support transition, the cellular
\(E_{2}\)-centre input, the Nygaard \(2\)-convergence input, and the
flat-locus centre input. The last three are stated here in exactly the
form required for the arithmetic transition locus.
\end{remark}

\section{Target G.2: BC-Diamond-Glue Final Closure}
\label{v4-arith:w6-g:bcglue-final}

\subsection{Chenevier--Bella\"iche Bernstein--Galois Decomposition}

\begin{theorem}[Bella\"iche--Chenevier 2009 Bernstein--Galois decomposition]
\label{v4-arith:w6-g:thm:bc-bernstein-galois}
\ClaimStatusProvedElsewhere. Bella\"iche--Chenevier 2009 Ast\'erisque
324, Chapter 8 (``Bernstein families and Galois
representations''), Thm.~8.1.1 proves: for each rigid point \(x\) of
the Chenevier pseudocharacter ring
\(\mathrm{Spec}(R^{\mathrm{Ch}})\) over \(W(k)\) with residual
determinant-trace pseudocharacter \(\bar{\mathcal{P}}\) absolutely
irreducible, and for each Bernstein block \(\mathfrak{s}\) of the
local Hecke algebra at a prime \(p\) in the finite set of ramified
primes of \(x\), there is a canonical direct-summand decomposition of
\(R^{\mathrm{Ch}}\)-modules
\begin{equation}
\label{v4-arith:w6-g:eq:bc-bernstein-galois}
R^{\mathrm{Ch}}\text{-}\mathrm{Mod}^{\mathrm{ss}}_{x}\;=\;\bigoplus_{\mathfrak{s}\in\mathfrak{B}_{p}(x)}R^{\mathrm{Ch}}_{\mathfrak{s}}\text{-}\mathrm{Mod}^{\mathrm{ss}}_{x},
\end{equation}
where \(R^{\mathrm{Ch}}_{\mathfrak{s}}\) is the Bernstein-block
component of \(R^{\mathrm{Ch}}\), locally presented by a
Galois--Bernstein quotient of the universal deformation ring.
Globally, this decomposition assembles to a product decomposition
\begin{equation}
\label{v4-arith:w6-g:eq:bc-global-bernstein-galois}
R^{\mathrm{Ch}}\;\simeq\;\prod_{\text{Bernstein data }\vec{\mathfrak{s}}}R^{\mathrm{Ch}}_{\vec{\mathfrak{s}}},
\end{equation}
the product over Bernstein data
\(\vec{\mathfrak{s}}=(\mathfrak{s}_{v})_{v}\) at every ramified place,
with \(R^{\mathrm{Ch}}_{\vec{\mathfrak{s}}}\) the corresponding
Galois--Bernstein component of the Chenevier pseudocharacter ring.
The product is restricted: for almost all \(v\), \(\mathfrak{s}_{v}\)
is forced to be the unramified-principal-series block.
\end{theorem}

\subsection{Francis--Gaitsgory Factorization Gluing}

\begin{theorem}[factorization gluing input for Bernstein centres]
\label{v4-arith:w6-g:thm:fg-factorization-gluing}
\ClaimStatusConjectured. Francis--Gaitsgory arXiv:1103.5803
Thm.~A (``factorization homology on manifolds is
\(E_{n}\)-monoidal'') combined with their Prop.~4.2 (``factorization
gluing is a limit along the factorization diagram on finite subsets'')
is used here in the following arithmetic form: for a
collection \(\{\mathcal{Z}_{v}\}_{v}\) of \(E_{2}\)-algebras in
presentable stable \(\infty\)-categories, the Francis--Gaitsgory
factorization product
\begin{equation}
\label{v4-arith:w6-g:eq:fg-factorization-product}
\int^{\mathrm{FG}}_{|\overline{\mathrm{Spec}\,\mathbb{Z}}|}\{\mathcal{Z}_{v}\}\;=\;\varinjlim_{S\text{ finite}}\Bigl(\bigotimes'_{v\in S}\mathcal{Z}_{v}\Bigr)
\end{equation}
is an \(E_{2}\)-algebra that glues the pointwise \(\mathcal{Z}_{v}\)
along the discrete-base factorization axiom. Pointwise compatibility
(for distinct \(v\neq w\), the factorization structure is the tensor
product \(\mathcal{Z}_{v}\otimes\mathcal{Z}_{w}\)) is automatic on a
discrete base; the filtered colimit over finite subsets \(S\) is the
factorization homology itself.

Applied with \(\mathcal{Z}_{v}\) the local Bernstein centre at each
finite place \(v\in|\overline{\mathrm{Spec}\,\mathbb{Z}}|\) and the
archimedean \(\mathcal{Z}_{\infty}=\mathbf{Z}(\mathcal{U}(\mathfrak{gl}_{2}))\)
(the Harish-Chandra centre), the factorization product is the
global Bernstein centre:
\begin{equation}
\label{v4-arith:w6-g:eq:fg-global-bernstein}
\mathcal{Z}^{\mathrm{FG}}_{\mathrm{glob}}\;=\;\int^{\mathrm{FG}}_{|\overline{\mathrm{Spec}\,\mathbb{Z}}|}\{\mathcal{Z}_{v}\}\;\simeq\;\varinjlim_{S\text{ finite}}\bigotimes'_{v\in S}\mathcal{Z}_{v}.
\end{equation}
\end{theorem}

\subsection{Gluing Across Prismatic Charts and Untilts}

\begin{theorem}[BC-diamond-glue across prismatic charts and untilts]
\label{v4-arith:w6-g:thm:bcglue-across-charts}
\ClaimStatusProvedHereConditional. Conditional on
\(\mathrm{L}_{\mathrm{ACD}}\).\(0\)--.\(3\) via
Chapter~\ref{v4-arith:w5-e:chapter}
(\Cref{v4-arith:w5-e:thm:LACD-cuspidal}), and on (AG.det.transition)
(\Cref{v4-arith:w6-g:cor:agdet-closed}), and on (BCglue.local)
and (BCglue.glue),
the two sides of (BC-diamond-glue.core)
(\Cref{v4-arith:bcglue-k5:thm:core-compatibility}) glue across
prismatic charts and Fargues--Fontaine untilts:
\begin{enumerate}
\item The global Bernstein centre
\(\mathcal{Z}_{\mathrm{glob}}\) acting on
\(\mathrm{D}^{\mathrm{sph\text{-}fin}}([GL_{2}])^{\mathrm{cusp}}\),
identified via Chenevier--Bella\"iche Bernstein--Galois
decomposition
(\Cref{v4-arith:w6-g:thm:bc-bernstein-galois}) with a
Bernstein-product decomposition of \(R^{\mathrm{Ch}}\);
\item The factorization-homology global centre
\(\mathcal{Z}^{\mathrm{FG}}_{\mathrm{glob}}\) of
\Cref{v4-arith:w6-g:thm:fg-factorization-gluing}, the filtered
colimit of restricted tensor products of local Bernstein centres;
\item The \(E_{2}\)-centre
\(\mathcal{Z}^{E_{2}}(\mathrm{IndCoh}_{\mathrm{Nilp}}(\mathcal{X}^{\mathrm{Sel}}_{\bar\rho,\det,\mathcal{L}}))\)
of the Selmer-derived global spectral stack, identified via
(AG.det.transition)
(\Cref{v4-arith:w6-g:cor:agdet-closed}) with the fibre-product
\(E_{2}\)-centre on the determinantal cleavage;
\end{enumerate}
all three \(E_{2}\)-algebras coincide:
\begin{equation}
\label{v4-arith:w6-g:eq:three-centres-coincide}
\mathcal{Z}_{\mathrm{glob}}\;\simeq\;\mathcal{Z}^{\mathrm{FG}}_{\mathrm{glob}}\;\simeq\;\mathcal{Z}^{E_{2}}\bigl(\mathrm{IndCoh}_{\mathrm{Nilp}}(\mathcal{X}^{\mathrm{Sel}}_{\bar\rho,\det,\mathcal{L}})\bigr).
\end{equation}
\end{theorem}

\begin{proof}
\emph{Step 1: \(\mathcal{Z}_{\mathrm{glob}}\simeq\mathcal{Z}^{\mathrm{FG}}_{\mathrm{glob}}\)
via L\(_{\mathrm{ACD}}\).} The three \(E_{2}\)-algebras
\(\mathcal{Z}_{\mathrm{glob}}\), \(\mathcal{Z}^{\mathrm{FG}}_{\mathrm{glob}}\),
and \(\mathcal{Z}^{E_{2}}(\mathrm{IndCoh}_{\mathrm{Nilp}})\) are
\(E_{2}\)-algebras in presentable stable \(\infty\)-categories;
their common underlying object is \(\mathcal{V}^{\mathrm{prim}}\),
the Hochschild trace class of the identity on the global arithmetic
Iwahori--Hecke category of \(GL_{2}\), which is not a vertex algebra
but a Hochschild trace class whose character is an \(L\)-function.
By
\Cref{v4-arith:w5-e:thm:LACD-cuspidal} (L\(_{\mathrm{ACD}}\) on the
cuspidal sph-fin locus), the restricted tensor product of local
sph-fin categories is equivalent to the cuspidal sph-fin automorphic
category, and the equivalence is compatible with the Bernstein-centre
action at each prime. Taking \(E_{2}\)-centres and using the
factorization-gluing formula
(\ref{v4-arith:w6-g:eq:fg-global-bernstein}), the global Bernstein
centre acting on
\(\mathrm{D}^{\mathrm{sph\text{-}fin}}([GL_{2}])^{\mathrm{cusp}}\)
agrees with the factorization centre
\(\mathcal{Z}^{\mathrm{FG}}_{\mathrm{glob}}\).

\emph{Step 2: \(\mathcal{Z}^{\mathrm{FG}}_{\mathrm{glob}}\simeq R^{\mathrm{Ch}}\)
on the cuspidal locus.} By Bella\"iche--Chenevier 2009 Chapter 8
Thm.~8.1.1
(\Cref{v4-arith:w6-g:thm:bc-bernstein-galois}), the Bernstein--Galois
decomposition
(\ref{v4-arith:w6-g:eq:bc-global-bernstein-galois}) identifies the
Chenevier pseudocharacter ring \(R^{\mathrm{Ch}}\) as a product of
Bernstein components indexed by Bernstein data. Each component is the
image of the tensor product of local Bernstein centres at the
ramified places; almost-all unramified places contribute the
spherical-unit subring. The factorization centre, by
(\ref{v4-arith:w6-g:eq:fg-factorization-product}), is the filtered
colimit of these tensor products; its image inside the cuspidal
automorphic centre coincides with the Bernstein-product decomposition
of \(R^{\mathrm{Ch}}\).

\emph{Step 3: \(R^{\mathrm{Ch}}\simeq\mathcal{Z}^{E_{2}}(\mathrm{IndCoh}_{\mathrm{Nilp}})\)
via (Ch1) and (AG.det.transition).} By (Ch1)
(\Cref{v4-arith:w5-h:thm:ch1-closure}), the cross-prime Chenevier
restriction is a flat monomorphism; by
(AG.det.transition)
(\Cref{v4-arith:w6-g:thm:agdet-transition-closure}), the
\(E_{2}\)-centre of \(\mathrm{IndCoh}_{\mathrm{Nilp}}\) on the
determinantal cleavage is the fibre product of the three centres.
The resulting fibre product maps to \(R^{\mathrm{Ch}}\) via the
Chenevier invariant projection
(\Cref{v4-arith:bcglue-core:thm:global-bernstein-chenevier}): the
\(GL_{2}\)-invariant ring of the Selmer-derived global stack is
\(R^{\mathrm{Ch}}\), and the \(E_{2}\)-centre restricts to \(R^{\mathrm{Ch}}\)
on the invariant sector. The map is an isomorphism by the
combined Bella\"iche--Chenevier Bernstein--Galois decomposition and
(Ch1) flatness.

\emph{Step 4: Prismatic-chart gluing.} The prismatic structure at
each prime \(p\) gives a separate chart of \(R^{\mathrm{Ch}}\)
(Chenevier 2014 + Bhatt--Lurie 2022). Cross-prime gluing of the
charts is the statement (Ch1), plus the prismatic-flatness
\(\Pi\)-flat hypothesis
(\Cref{v4-arith:bcglue-k5:lem:prismatic-flatness}). Under these two
inputs, the charts glue via the Chenevier pseudocharacter
restriction \(R^{\mathrm{Ch}}\hookrightarrow R^{\mathrm{Ch}}_{p}\otimes R^{\mathrm{Ch}}_{q}\)
(\Cref{v4-arith:w5-h:thm:ch1-closure} and
\Cref{v4-arith:bcglue-core:thm:prismatic-nygaard-transition}).

\emph{Step 5: Fargues--Fontaine untilt gluing.} The
Fargues--Fontaine curve \(X_{\mathrm{FF},p}\) at \(p\) provides the
untilt geometry over \((\mathbb{Z}_{p},(p))\); gluing across untilts is
a diamond-level gluing, supplied by Fargues--Scholze arXiv:2102.13459
Thm.~II.2.19 (diamond gluing along \(\mathrm{Spd}(\mathbb{Z}_{p})\)).
On the centre side, the untilt gluing is trivial on the
\(GL_{2}\)-invariant sector: the Chenevier pseudocharacter ring is
invariant of the untilt choice (Chenevier 2014 Rem.~4.3).

\emph{Step 6: Combined chain.} Combining Steps 1--5:
\begin{align*}
\mathcal{Z}_{\mathrm{glob}}
&\xrightarrow{\text{L}_{\mathrm{ACD}}\text{ (Step 1)}}
\mathcal{Z}^{\mathrm{FG}}_{\mathrm{glob}}\\
&\xrightarrow{\text{BC Thm.~8.1.1 (Step 2)}}
R^{\mathrm{Ch}}\\
&\xrightarrow{\text{(Ch1)}+\text{(AG.det.transition) (Step 3)}}
\mathcal{Z}^{E_{2}}\bigl(\mathrm{IndCoh}_{\mathrm{Nilp}}(\mathcal{X}^{\mathrm{Sel}}_{\bar\rho,\det,\mathcal{L}})\bigr),
\end{align*}
each arrow an isomorphism under the stated conditional inputs. The
three-way coincidence
(\ref{v4-arith:w6-g:eq:three-centres-coincide}) follows.

\emph{Independence.} Proof source: Bella\"iche--Chenevier
2009 Chapter 8 (Bernstein--Galois decomposition via Galois
cohomology of pseudocharacter families). Verification sources:
Francis--Gaitsgory arXiv:1103.5803 (factorization homology on
manifolds; no Galois input) and (Ch1)
(\Cref{v4-arith:w5-h:thm:ch1-closure}) + (AG.det.transition)
(\Cref{v4-arith:w6-g:thm:agdet-transition-closure})
(both at the \(E_{2}\)-centre level with independent source lists).
The three lists are disjoint: Galois cohomology of families for
BC; factorization homology for FG; prismatic/cellular \(E_{k}\) for the
inherited centre comparisons.
\end{proof}

\begin{corollary}[(BCglue.core) from the standing inputs]
\label{v4-arith:w6-g:cor:bcglue-core-closed}
\ClaimStatusProvedHereConditional. Under the hypotheses of
\Cref{v4-arith:w6-g:thm:bcglue-across-charts}
(L\(_{\mathrm{ACD}}\).\(0\)--.\(3\), (AG.det.transition),
(BCglue.local) and (BCglue.glue) stood on their conditional inputs,
and arithmetic AG at the centre), the core compatibility
(BCglue.core) (\Cref{v4-arith:bcglue-core:cj:core}) follows.
\end{corollary}

\begin{proof}
By \Cref{v4-arith:bcglue-core:thm:core-final-reduction}, (BCglue.core)
is equivalent to (Ch1) \(\wedge\) (AG.det). The first is
\Cref{v4-arith:w5-h:thm:ch1-closure}; the second is
\Cref{v4-arith:w6-g:cor:agdet-closed}. Hence (BCglue.core)
follows, which by \Cref{v4-arith:w6-g:thm:bcglue-across-charts} is
the three-way centre coincidence
(\ref{v4-arith:w6-g:eq:three-centres-coincide}).
\end{proof}

\begin{corollary}[BC-diamond-glue from three conditional inputs]
\label{v4-arith:w6-g:cor:bcglue-final}
\ClaimStatusConditional. Under the three conditional inputs
(BCglue.local) (Conditional), (BCglue.glue) (Conjectural,
conditional on \(\Pi\)-flat), and arithmetic AG at the centre
(Conjectural), the BC-diamond-glue compatibility
(\Cref{v4-arith:acd:cj:bc-diamond-glue-concrete}) follows from
\Cref{v4-arith:w6-g:cor:agdet-closed} and
\Cref{v4-arith:w6-g:thm:bcglue-across-charts}.
\end{corollary}

\begin{proof}
Assemble \Cref{v4-arith:bcglue-k5:thm:decomposition} (the three-part
decomposition of BC-diamond-glue) with the reductions of the three
parts: (BCglue.local) and (BCglue.glue) at their conditional stage
from \Cref{v4-arith:bcglue-k5:thm:decomposition}; (BCglue.core) via
\Cref{v4-arith:w6-g:cor:bcglue-core-closed}.
\end{proof}

\section{Conditional Locus}
\label{v4-arith:w6-g:residuals}

\begin{theorem}[\(\mathrm{L}_{\mathrm{ACD}}\) conditional locus]
\label{v4-arith:w6-g:thm:lacd-wave6}
\ClaimStatusConjectured (conditional on the named standing residuals).
\(\mathrm{L}_{\mathrm{ACD}}\) has the following conditional locus:
\begin{enumerate}
\item (BCglue.local):
\Cref{v4-arith:bcglue-k5:thm:local-BZCHN-plus-EG}.
\item (BCglue.glue): Conjectural, conditional on \(\Pi\)-flat
\Cref{v4-arith:bcglue-k5:thm:prismatic-gluing}).
\item (Ch1): Conditional on the pseudocharacter trace comparison
(\Cref{v4-arith:w5-h:thm:ch1-closure}).
\item (AG.det): Conditional on \(\Pi\)-flat, (BCglue.local),
and arithmetic AG at the centre
(\Cref{v4-arith:w6-g:cor:agdet-closed}).
\item F2.c:
(\Cref{v4-arith:tate:thm:F2c-main} /
\Cref{v4-arith:w5-h:thm:f2c-e2-transport}).
\item L\(_{\mathrm{ACD}}\).\(0\)--.\(3\): On the cuspidal sph-fin
locus (\Cref{v4-arith:w5-e:thm:LACD-cuspidal}); Eisenstein extension
is \Cref{v4-arith:w5-e:conj:LACD-Eisenstein}.
\item BC-diamond-glue: Conditional on
(BCglue.local), (BCglue.glue), and arithmetic AG at the centre
(\Cref{v4-arith:w6-g:cor:bcglue-final}).
\end{enumerate}
\end{theorem}

\begin{proof}
Assemble \Cref{v4-arith:w5-h:cor:LACD-wave5} with
\Cref{v4-arith:w6-g:cor:agdet-closed},
\Cref{v4-arith:w6-g:cor:bcglue-core-closed}, and
\Cref{v4-arith:w6-g:cor:bcglue-final}.
\end{proof}

\begin{remark}[remaining inputs]
\label{v4-arith:w6-g:rem:domain-delta}
(BCglue.local), (BCglue.glue), and arithmetic AG at the centre are the
three independent inputs in the BC-diamond-glue comparison. The
determinant transition supplies the codimension-two \(E_{2}\)-centre
compatibility once the cellular, Nygaard, and flat-locus centre inputs
are imposed.
\end{remark}

\begin{remark}[arithmetic and Eisenstein inputs]
\label{v4-arith:w6-g:rem:what-wave6-cannot-touch}
(BCglue.local) and (BCglue.glue) remain as conditional inputs:
these are primary-source transition theorems in the local-Langlands and
prismatic-cohomology literatures. Arithmetic AG at the
centre is a standing conditional input; the
three-way centre coincidence
(\ref{v4-arith:w6-g:eq:three-centres-coincide}) is conditional on
it through Step 1 (L\(_{\mathrm{ACD}}\)) and Step 3
((AG.det.transition) via (AG.det) reduction).

The Eisenstein extension of
\(\mathrm{L}_{\mathrm{ACD}}\)
(\Cref{v4-arith:w5-e:conj:LACD-Eisenstein}) is a
Langlands--Shahidi question about intertwining operators on
continuous-spectrum automorphic representations, requiring a
separate argument.
\end{remark}

\section{Independent Source Data}
\label{v4-arith:w6-g:data}

\begin{remark}[independent source data]
\label{v4-arith:w6-g:rem:data}
The independent source data consist of three separations:
\begin{itemize}
\item For \Cref{v4-arith:w6-g:thm:agdet-transition-closure}, the
singular-support source is Arinkin--Gaitsgory arXiv:1201.6343 \S 11,
and the comparison sources are Galatius--Kupers--Randal-Williams
arXiv:1805.07184, Bhatt--Lurie arXiv:2201.06120, and
Bezrukavnikov--Finkelberg--Mirkovi\'c 2005.
\item For \Cref{v4-arith:w6-g:thm:bcglue-across-charts}, the
Bernstein--Galois source is Bella\"iche--Chenevier 2009 Chapter 8, and
the comparison source is Francis--Gaitsgory arXiv:1103.5803 together
with the centre comparisons above.
\item For \Cref{v4-arith:w6-g:cor:bcglue-final}, the independent
conditional inputs are (BCglue.local), (BCglue.glue), and arithmetic AG
at the centre.
\end{itemize}
\end{remark}

\section{Bibliography Stanza}
\label{v4-arith:w6-g:bibliography}

The primary literature invoked in this chapter:

\begin{description}
\item[Arinkin, D.\ and Gaitsgory, D.] (2015). \emph{Singular support
of coherent sheaves and the geometric Langlands conjecture}. Selecta
Math. 21, 1--199. arXiv:1201.6343. \S 10 (nilpotent-support
decomposition), \S 11 (boundary structure at codim 2), \S 11.4.1
(Kostant section) used in
Section~\ref{v4-arith:w6-g:agdet-transition}.
\item[Galatius, S., Kupers, A., and Randal-Williams, O.] (2024).
\emph{Cellular \(E_{k}\)-algebras}. Ann. Math. 200, arXiv:1805.07184.
Thm.~17.3 (\(E_{2}\)-centre via cellular \(E_{k}\)-algebras), \S 11.3
(\(E_{k}\)-cellular convergence), \S 18.4 (higher-codim cellular
\(E_{k}\)-resolutions) used in
\Cref{v4-arith:w6-g:thm:gkrw-2-convergence}.
\item[Bhatt, B.\ and Lurie, J.] (2022). \emph{Absolute prismatic
cohomology}. arXiv:2201.06120. \S 8.3--8.5 (Nygaard \(2\)-convergence),
Cor.~8.9 (spectral-sequence abutment) used in
\Cref{v4-arith:w6-g:thm:bl-2-convergence}.
\item[Bezrukavnikov, R., Finkelberg, M., and Mirkovi\'c, I.] (2005).
\emph{Equivariant homology and \(K\)-theory of affine Grassmannians and
Toda lattices}. Compositio Math. 141. Thm.~1.1
(equivariant cohomology = centre on the flat locus) used in
\Cref{v4-arith:w6-g:thm:bfm-equivariant}.
\item[Bella\"iche, J.\ and Chenevier, G.] (2009). \emph{Families of
Galois representations and Selmer groups}. Ast\'erisque 324.
Chapter 8 Thm.~8.1.1 (Bernstein--Galois decomposition of
pseudocharacter families) used in
\Cref{v4-arith:w6-g:thm:bc-bernstein-galois}.
\item[Francis, J.\ and Gaitsgory, D.] (2012). \emph{Chiral Koszul
duality}. Selecta Math. 18, arXiv:1103.5803. Thm.~A (factorization
homology on manifolds), Prop.~4.2 (factorization gluing as limit on
finite subsets) used in
\Cref{v4-arith:w6-g:thm:fg-factorization-gluing}.
\item[Chenevier, G.] (2014). \emph{The \(p\)-adic analytic space of
pseudo-characters of a profinite group}. JAMS 27, 1007--1074.
Chenevier 2014 enters through
\Cref{v4-arith:w5-h:thm:ch1-closure}.
\item[Fargues, L.\ and Scholze, P.] (2021). \emph{Geometrization of
the local Langlands correspondence}. arXiv:2102.13459. Thm.~II.2.19
(diamond gluing along \(\mathrm{Spd}(\mathbb{Z}_{p})\)) used in
\Cref{v4-arith:w6-g:thm:bcglue-across-charts} Step~5.
\item[Zhu, X.] (2017). \emph{Affine Grassmannians and the geometric
Satake in mixed characteristic}. Ann. Sci.\ ENS 50, arXiv:1407.8519.
\(p\)-adic geometric Satake used in
\Cref{v4-arith:w6-g:thm:bfm-equivariant} for the arithmetic analog
of the flat-locus embedding.
\item[Lurie, J.] (2017). \emph{Higher algebra}. 5.3.1.15
(fibre-product commutes with \(\mathcal{Z}^{E_{2}}\)) used in
\Cref{v4-arith:w6-g:thm:agdet-transition-closure} Step~5.
\end{description}

The seven independent lists
(AG, GKRW, BL, BFM, BC, FG, Lurie) have no primary-author overlap,
no derivation dependency between them, and underlie formalisms
that were built independently (singular support on derived stacks;
cellular \(E_{k}\)-algebras; prismatic cohomology; affine-Grassmannian
equivariant homology; Galois pseudocharacter families; factorization
homology on manifolds; \(\infty\)-algebra of operads). No citation
appears at more than one of the (AG.det.transition) and
(BC-diamond-glue final) source pairs.

\section{Consequences}
\label{v4-arith:w6-g:summary}

\begin{enumerate}
\item (AG.det.transition) follows from the four-source assembly
AG 2015 + the cellular input + the Nygaard input + the flat-locus input
on the codim-\(\leq 2\)
transition locus of the determinantal cleavage
(\Cref{v4-arith:w6-g:thm:agdet-transition-closure}). This gives
(AG.det) conditional on \(\Pi\)-flat and (BCglue.local)
(\Cref{v4-arith:w6-g:cor:agdet-closed}).
\item BC-diamond-glue follows from the three-way centre
coincidence
(\ref{v4-arith:w6-g:eq:three-centres-coincide})
through Bella\"iche--Chenevier Bernstein--Galois decomposition +
Francis--Gaitsgory factorization gluing + inherited (Ch1) +
(AG.det.transition) (\Cref{v4-arith:w6-g:thm:bcglue-across-charts}).
This gives (BCglue.core) under the standing inputs
(\Cref{v4-arith:w6-g:cor:bcglue-core-closed}) and hence the full
BC-diamond-glue comparison (\Cref{v4-arith:w6-g:cor:bcglue-final}).
\item The remaining independent inputs are (BCglue.local),
(BCglue.glue), and arithmetic AG at the centre.
\item The source separation is documented in
Section~\ref{v4-arith:w6-g:bibliography}.
\end{enumerate}
